% =================== Kullanılan Notasyonlar ===================
\chapter*{Kullanılan Notasyonlar}
\addcontentsline{toc}{chapter}{Kullanılan Notasyonlar}

Bu çalışmada yerçekimi alanlarının küresel harmonikler kullanılarak modellenmesi kapsamında aşağıda verilen notasyonlar kullanılmıştır. Semboller, metin boyunca tutarlı olacak şekilde tanımlanmış olup, aksi belirtilmedikçe skaler büyüklükleri temsil etmektedir.

\medskip
Özellikle Legendre polinomları ve ilişkili fonksiyonlarda, fonksiyonel bağımlılığın cebirsel çarpma işlemleriyle karıştırılmasını önlemek ve matematiksel okunabilirliği
artırmak amacıyla, fonksiyon argümanları standart parantez $()$ yerine köşeli parantez $[]$ ile (örneğin $P_{\ell}^{m}[\sin\phi_{gc}]$) gösterilmiştir.

\vspace{0.4cm}

\setlength{\LTpre}{0pt}
\setlength{\LTpost}{0pt}
\renewcommand{\arraystretch}{1.15}

\begin{longtable}{>{\raggedright\arraybackslash}p{3.2cm} p{10.8cm}}
\toprule
\textbf{Sembol} & \textbf{Açıklama} \\
\midrule
\endfirsthead

\toprule
\textbf{Sembol} & \textbf{Açıklama} \\
\midrule
\endhead

\midrule
\multicolumn{2}{r}{\emph{(devam ediyor)}}\\
\endfoot

\bottomrule
\endlastfoot

$r$ & Cisim merkezinden ölçülen radyal uzaklık (skaler) \\
$\mathbf{r}$ & Konum vektörü \\

$R_E$ & Dünya için referans yarıçapı \\
$R_M$ & Ay için referans yarıçapı \\
$R_{\text{ref}}$ & Genel referans yarıçapı (model bazlı; Dünya/Ay) \\

$\phi_{gc}$ & Geosentrik enlem \\
$\lambda$ & Boylam \\
$\theta$ & Kolatitude ($\theta=\tfrac{\pi}{2}-\phi_{gc}$) \\
$x$ & Legendre argümanı ($x=\cos\theta=\sin\phi_{gc}$) \\

$U$ & Yerçekimi potansiyeli \\
$\mu$ & Standart çekim parametresi ($\mu = GM$) \\
$G$ & Evrensel çekim sabiti \\
$M$ & Gökcisminin kütlesi \\

$P_{\ell}[x]$ & $\ell$ dereceli Legendre polinomu \\
$P_{\ell}^{m}[x]$ &
$\ell$ derece, $m$ mertebeli ilişkili Legendre fonksiyonu (Condon--Shortley fazı hariç) \\
$P_{\ell}^{m,\mathrm{CS}}[x]$ &
Condon--Shortley fazını içeren tanım ($P_{\ell}^{m,\mathrm{CS}}[x]=(-1)^m P_{\ell}^{m}[x]$) \\

$\overline{P}_{\ell}^{m}[x]$ &
Tam-normalize edilmiş ilişkili Legendre fonksiyonu (real harmonikler) \\
$N_{\ell,m}$ & Normalizasyon çarpanı ($\overline{P}_{\ell}^{m}=N_{\ell,m}P_{\ell}^{m}$) \\

$\ell$ & Küresel harmonik derecesi (degree) \\
$m$ & Küresel harmonik mertebesi (order) \\

$C_{\ell,m}$ & Kosinüs küresel harmonik katsayısı (normalize edilmemiş baz ile uyumlu) \\
$S_{\ell,m}$ & Sinüs küresel harmonik katsayısı (normalize edilmemiş baz ile uyumlu) \\
$\overline{C}_{\ell,m}$ & Tam-normalize edilmiş kosinüs katsayısı (EGM/jeodezik konvansiyon) \\
$\overline{S}_{\ell,m}$ & Tam-normalize edilmiş sinüs katsayısı (EGM/jeodezik konvansiyon) \\

$J_{\ell}$ & Zonal harmonik katsayısı ($m=0$ terimleri ile ilişkili) \\

$d_Q$ & Gözlem noktası ile kütle elemanı arasındaki mesafe (skaler) \\

$\mathbf{a}$ & Yerçekimi ivme vektörü \\

\end{longtable}

\vspace{0.2cm}

Bu notasyonlar, Dünya ve Ay için kullanılan tüm gravite modelleri ve yörünge
analizlerinde ortak bir referans çerçevesi sağlamak amacıyla tanımlanmıştır.
